\documentclass[a4paper,12pt]{ctexart}
\usepackage{xcolor}
\usepackage{graphicx}
\usepackage[top=1.8cm, bottom=1.8cm, left=1.8cm, right=1.8cm]{geometry}
\usepackage{array}
\usepackage{hyperref}
\hypersetup{colorlinks=true,linkcolor=blue!70!black}
\linespread{1.35}

\newcommand{\ruby}[2]{#1\textsuperscript{#2}}

\newcommand{\shihai}[2]{%
  \vspace{0.35cm}
  \noindent
  \fcolorbox{orange!60!black}{orange!8}{%
    \begin{minipage}{0.915\textwidth}
    \textbf{\textcolor{orange!70!black}{史海拾遗：}#1} \\[0.1cm]
    {\small #2}
    \end{minipage}%
  }
  \vspace{0.35cm}
}

\title{\textcolor{blue!70!black}{\Huge\textbf{政治与法律入门}}}
\author{}
\date{}

\begin{document}
\maketitle
\thispagestyle{empty}

\tableofcontents
\newpage

\section{\textcolor{orange!80!black}{第一课\quad 为什么要有政府？——从自然状态讲起}}

\subsection*{一个没有政府的世界}

想象一个世界：没有警察、没有法院、没有交通管理局。你家门口的马路坏了——没人修。你隔壁的人半夜用冲击钻拆墙——没人管。远处的工厂往你家门口的河里直接倾倒废酸——你能怎么办？你一个人去找工厂的保安理论——他把你轰出去。

17世纪的英国哲学家\textbf{托马斯·霍布斯}把这个假想状态叫做"\textbf{自然状态}"。他的描述很不乐观：在没有一个共同权力来使所有人畏惧服从的地方，人与人之间会处在"\textbf{所有人对所有人的战争}"状态。生活将是"\textbf{孤独、贫困、污秽、野蛮和短暂的}"。

\subsection*{一个真实的"自然状态"——索马里}

霍布斯写这些的时候，真正持久的"自然状态"在地球上很难找到。但现代历史提供过一个严重接近它的例子：\textbf{1991年索马里}。西亚德·巴雷政权被推翻之后，这个东非国家进入了一段没有统一中央政府的时间。各个地方军阀割据，同一个港口可能有两三个"海关"分别拦截货车并向同一车粮食征收几遍税费。因为没有统一的司法系统，商人和居民只好求助于\textbf{氏族长老}用传统习惯法来调解土地和血仇纠纷。没有中央银行发行货币——市场上同时流通着几版早已超过使用年限的纸币，以及战前铸的铜币。\textbf{索马里这几十年的经历是一个反面的案例：当一个国家长期没有统一的可执行规则的政府时，哪怕基础贸易都会退回到打火机、子弹、骆驼这些最最原始的以物易物或武装护送商队的形式。}

\subsection*{从自然状态到社会契约}

霍布斯的解方很直接：\textbf{把所有人的暴力权集中交给一个人或一个会议——利维坦（一个用强制力保障所有人生命安全的庞然大物）}。作为交换，个人不再私自动用暴力解决问题。但这个交出"全部权利"的契约只保了安全，不一定保自由。\textbf{在霍布斯的契约里——国王要你干什么——你不服从就可能死。}

几十年后另一位英国哲学家\textbf{约翰·洛克}修正了这张合同书。洛克认为自然状态不是混乱至极的杀人场——而是一块缺少了"公认裁判员"的绿地。人们为了保护自己的\textbf{生命、自由和财产}，自愿把一部分权利转交给政府。\textbf{但政府违反了对人民的信托时——人民有权利把它换掉。}

再隔了一个世纪，法国思想家\textbf{卢梭}把合同的逻辑翻了一面：\textbf{真正的主权者是全体人民组成的"公意"}。政府不过是受委托的执行人。卢梭不认为人们应该交出权利给一个国王——他认为权利应该始终停留在集体手里。\textbf{霍布斯要秩序、洛克要自由、卢梭要人民自己做主。这三份契约书的争吵，到今天还在世界各国的宪法条文和议会辩论里回响。}

\xuehai{社会契约——一种假想的"同意"}{
  社会契约论不是历史记录——没有人真的签过一份叫"和政府签的合同"的文件。它是政治哲学里用来回答一个根本问题的一套逻辑工具：如果权力不是来自上帝（君权神授），也不该来自纯粹的武力——那一个政府凭什么要求我服从？

  契约论者的共同策略是：假设我们回到政府成立之前的状态——然后推理所有理性的人会同意怎样一套规则。霍布斯说"把所有暴力交出去换安全"。洛克说"交一部分——政府违约人民可以换掉它"。卢梭说"权利不应该交给任何人——应该由全体公民共同掌握"。三份合同，同一套推理装置，三个完全不同的政治设计。

  社会契约是一个虚构的工具——但它在真实历史上产生了真实的爆炸力。美国1776年的《独立宣言》和法国1789年的《人权宣言》都用——甚至可以说直接抄了——洛克和卢梭的修辞。一份假想的合同，被墨水印在真实的羊皮纸上，然后被枪炮和断头台一并捍卫。
}

\subsection*{要点总结}
\begin{itemize}
  \item 霍布斯设想的"自然状态"是所有人对所有人的战争——因此需要一个强有力的"利维坦"来维持秩序
  \item 洛克修正了霍布斯：人们交出部分权利以保护生命、自由和财产——但政府违约时人民有权更换它
  \item 卢梭进一步：主权属于全体人民的"公意"，政府只是受委托的执行人
  \item 三位思想家的共同前提：政府的合法性来自被统治者的同意（社会契约），而非神授或武力
  \item 索马里的案例说明：没有统一政府时，弱势者失去的是唯一可求助的第三方规则执行者
\end{itemize}

\subsection*{讨论题}
\begin{enumerate}
  \item 霍布斯要秩序、洛克要自由、卢梭要人民做主——这三种"契约精神"你更倾向哪一种？如果你能设计一个全新的国家，你会综合他们的哪些主张？
  \item 索马里没有政府时，商人求助于氏族长老用传统习惯法调解纠纷。这是不是说明人类社会中总会以某种形式出现"规则和权威"？没有政府的"自由"和没有保护的"混乱"之间的线在哪里？
\end{enumerate}

\newpage
\section{\textcolor{orange!80!black}{第二课\quad 权力——谁统治、凭什么统治？}}

\subsection*{三种让你"服"的方式}

二十世纪初的德国社会学家\textbf{马克斯·韦伯}问了一个极其冷静的问题：你在马路上遇见一个穿制服的警察。你停下来。你为什么服从这个人的指令？不是因为他比你高大。是因为你相信他\textbf{有权}指挥。这个"相信"从哪里来的？

韦伯认为人类社会在不同的时代依赖三种完全不同的合法性依据：

\textbf{传统型：}"从来没有变过。所以是对的。"在位者的合法性靠的是古老的习俗链条——皇帝的儿子继位、部落酋长的后裔延续掌权——没人质问"你为什么有资格"，因为历史惯性把统治变成了习惯。这类型的弱点是\textbf{它很难自我革新}——当外部条件剧变——被入侵、饥荒、战争——传统权威不能迅速调整它的政策节奏，常常随之垮塌。

\textbf{魅力型：}"这个人就是不一样。"靠的是领袖本人——他的演讲、战功、人格。拿破仑回到法国的第一天报纸标题从"科西嘉怪物在儒昂港登陆"一路变成了"皇帝陛下即将抵达巴黎"——同一支笔，跟随态度的转向。魅力型的最大瓶颈是——\textbf{权力无法制度化}。当这个人消失——他的手下没有谁能复制他那套独一无二的光环。

\textbf{法理型：}"我是按大家事先同意的程序坐在这个位置上的。"\:你服从的是\textbf{椅子}——不是你椅子上那个具体的人。程序包括了选举、任命、考核、任期限制。如果坐在椅子上的人做得不好——同一套程序可以合法地把他换下去。\:\textbf{现代国家普遍建立在法理型的基座上——不管多党制还是一党领导——合格官员的任命和退出都需要有一条事先写清的程序。}

\xuehai{魅力型权威的"常规化"——从一个人到一套制度}{
  韦伯指出魅力型权威有一个内在的致命矛盾：它必须制度化才能延续——但制度化本身会杀死魅力。拿破仑靠个人战功和形象统治——他死后侄子路易·拿破仑只能靠模仿他的姿势和修辞来蹭残余的光环。当追随者无法回答"下一任魅力从哪里来"时——要么走向内战（亚历山大死后、成吉思汗死后），要么把魅力\textbf{转写到制度上}——"不是这个人本身了不起，是这个职位神圣。"法理的牙齿就是这条转化链条：把一个人的光芒磨成一行条款的签字权。
}

\subsection*{权力交接是最危险的时刻}

任何政体——在更替掌权者的时候——都像在飞行中拆发动机。世袭制的痛点：老国王三个儿子，下一个该传给谁？另外两个会不会带兵打回来？中国两千年的王朝史每写三页就会撞到一次因继承人模糊而炸开的藩王叛乱或外戚专权。法理型的难题同样锋利：\textbf{大选结果出来之后——输的那一方承不承认？}2021年1月美国国会山的冲击事件把这个问题从教科书拽到了全世界直播的屏幕上——承认败选并平稳交出行政档案不是一个自动发生的魔术——它需要每一个环节\u2014\u2014从选举管理机构、法院到军队——都默认遵守\textbf{同一条被写下来的规则}。

\subsection*{要点总结}
\begin{itemize}
  \item 韦伯提出三种权力合法性来源：传统型（习惯使然）、魅力型（个人魅力）、法理型（程序授权）
  \item 现代国家普遍建立在法理型权威基础上——不论谁坐在椅子上，权力来自程序而非个人
  \item 魅力型权威有一个致命矛盾：它必须制度化才能延续，但制度化本身会抹去"魅力"
  \item 权力交接是所有政体最危险的时刻——如何和平转移权力是一套制度成熟与否的试金石
\end{itemize}

\subsection*{讨论题}
\begin{enumerate}
  \item 你觉得一个"有魅力的独裁者"和一套"无趣但可靠的制度"哪个更能带来好政府？历史上有没有"魅力型领袖 + 法理型制度"的成功结合？
  \item 2021年美国国会山冲击事件中，输掉选举的一方不承认结果——你觉得是什么让一个政治体制能够在"输掉选举"的情况下仍然维持稳定？如果输的一方不认输，民主还能运转吗？
\end{enumerate}

\newpage
\section{\textcolor{orange!80!black}{第三课\quad 宪法——一个国家最根本的"源代码"}}

\subsection*{写下来的和没写下来的}

英国——至今没有一本叫"宪法"的单独成册文件。它的宪法内容散落在数百年来累积的\textbf{成文法（如大宪章1215年、权利法案1689年）}、法院判决先例和大量不成文惯例（如首相必须获得下议院信任才可执政，并没有法典明文规定）。宪法的灵活性在变革时代有好处：20世纪的福利国家制度、加入欧洲共同体、退出欧盟——每一项都是通过常规立法程序或全民公决实现的，没有需要触发"制宪大会"走繁复的修改路径。代价同样是明显的——如果某一天出现一个极度想集中权力的首相\u2014\u2014\textbf{"惯例"能不能挡住他？}

美国的制宪者在1787年写了一份极其短的文件——加上全部修正案全文不过几千词。但这几千词回答了几乎所有现代国家的根本难题：联邦政府管什么\u2014\u2014各州保留什么；立法、行政和司法三个机构互相重叠又互相咬合——任命官员、通过预算法案、审查法律是否违宪\u2014\u2014每一个动作被设计成了另一支齿轮的制动。这套机制的弱点\u2014\u2014三权分支在尖锐对立时会把政府的日常运转切成僵局（美国政府历史上多次因预算分歧而短暂"停摆"）。

\subsection*{宪法里的两条核心防线}

翻看任何当代宪法的条文，一定会找到对个人基本权利的明确列举——\textbf{不被任意逮捕、不被刑讯、表达不同意见的自由、信仰自由}。宪法和普通法律的分界线就在于：\textbf{普通法律可以被国会一纸表决改掉；宪法的基本权利条款要修改通常需要极高的超多数表决门槛（例如三分之二以上多数、或全国范围的公决）。}这层"硬壳"在保护谁？保护群体换了一届选票就无法被顺利欺负的那些人。

\subsection*{违宪审查——谁来管立法者自己？}

国会发布了一条法律——凡是政府认为"可能威胁公共安全"的人可以不经审判被关押最长五年。一个普通公民把这条法律告到最高法院。大法官翻开宪法的文本——"任何人的生命、自由和财产非经正当法律程序不得被剥夺。"判决是：\textbf{国会的这条法律违宪——作废。}违宪审查的本意不是"法官比议员聪明"。它是设置一道制动：\textbf{多数人投票通过的法案，也要受一部更高阶的文件的检查。}不同国家的审查者不同：德国是独立的宪法法院专审违宪案件；法国是宪法委员会；美国是普通联邦最高法院在审理具体案件中附带审查。\textbf{但背后的逻辑是共同的——民主不意味着立法机构可以做一切。}

\subsection*{要点总结}
\begin{itemize}
  \item 宪法是一个国家最根本的"源代码"——它规定了权力的分配、运行的规则和公民基本权利的保护底线
  \item 英国没有单一成文宪法文件，其宪法散见于成文法、判例和惯例中——灵活性高但依赖政治传统
  \item 美国宪法以三权分立为核心：立法、行政、司法三个机构互相制衡
  \item 宪法的基本权利条款修改门槛极高——这层"硬壳"保护了少数群体不被多数投票轻易剥夺权利
  \item 违宪审查制度允许司法机关否决立法机关通过的法律——这是"反多数"的制度设计，但目的是保护宪法底线
\end{itemize}

\subsection*{讨论题}
\begin{enumerate}
  \item 英国没有成文宪法却运转了几百年——你觉得对于一个国家来说，是"把规则全写下来"更重要，还是"有遵守规则的传统"更重要？如果传统和成文规则冲突了，哪个更可靠？
  \item 违宪审查是非民选的法官否决民选的议会——这个设计的合理性在于"保护少数人的基本权利不受多数人侵犯"。但谁来监督法官？如果法官自己违背了宪法，谁来纠正？
\end{enumerate}

\subsection*{中国的宪法框架——权力怎么分工的？}

1982年通过的现行《中华人民共和国宪法》也回答着和英美宪法同样的问题——权力怎么分配、怎么被约束——但它的结构不同。

全国最高的权力机关是\textbf{全国人民代表大会}。它负责制定和修改基本法律、审议政府工作报告、审查和批准国家预算与国民经济发展计划。\textbf{国务院}是最高国家行政机关——总理、各部委首长由全国人大决定或任命，向全国人大负责并报告工作。\textbf{国家主席}根据全国人大及其常委会的决定行使职权——公布法律、任免国务院组成人员、授予国家勋章、宣布进入紧急状态等。\textbf{中央军事委员会}领导全国武装力量。2018年修宪后增设了\textbf{国家监察委员会}——统一对所有行使公权力的公职人员进行监督、调查和处置，和原来的行政监察和反贪系统相比，监察委的覆盖范围扩展到了所有公权力的末端。

\textbf{最高人民法院}依据法律独立行使审判权，不受行政机关、社会团体和个人的干涉。\textbf{最高人民检察院}是国家的法律监督机关——对刑事案件审查批捕、提起公诉，对法院的生效判决和裁定进行监督，发现确有错误的可以依法提出抗诉。

这套架构的核心逻辑和"三权分立"不同——\textbf{它不强调三权之间的相互否决，而更侧重所有权力最终来源于全国人民代表大会，由其统一行使最高权力，再由各个具体机关在人大之下分工执行。}\:权力分在不同的机构之间，每一项权力的源头都必须在宪法条文上找到自己的出生证。

\newpage
\section{\textcolor{orange!80!black}{第四课\quad 法律是怎么工作的？——从立法到判决}}

\subsection*{一部法律怎么从一张纸变成有约束力的规则}

假设某市PM2.5年均值连续超标，市民和环保组织提请市人大制定更严格的工业排放控制。建议提交给城乡建设与环境资源保护专门委员会做立法前的实地调研取证。专家组论证草案之后形成征求意见稿向社会公开。公开渠道收到的几百上千条企业、居民和专家的反馈被编入修改建议表——经法制工作机构仔细筛选后交人大常委会逐条审议。

当常委会表决通过之后——这部法规会被正式公告指定生效日期。\textbf{一部地方性法规从最初始的民间建议到真正对工厂管用这个全过程往往贯穿了两三年，其间每一步的目的全是为了防止"拍脑袋立法"。}

\subsection*{民事诉讼、刑事诉讼和行政诉讼——三种法院里的角色}

\textbf{民事诉讼——"你和你的邻居因为倒垃圾而闹到法院"}。原告（你）诉被告（邻居），要求对方赔偿或恢复原状。胜诉标志可能是一纸判决书勒令被告为损坏的你的财产支付修复金额。民事案子里国家不是原告——它是裁判员。

\textbf{刑事诉讼——"检察院告张三入室盗窃"。}这不再是你和张三之间的私人纠纷。张三把邻居家的锁砸碎了，社会维度上每个人晚上锁门时的安全感都被他的行为多填了一层怀疑。国家（检察院代表国家）作为公诉方入场——立案侦查、批捕、预审、提起公诉、法院审判——每个节点都在规范张三的基本权利（不能虐待、不能超标超期羁押、有辩护权）。

\textbf{行政诉讼——"你觉得某个行政行为错了——你告作出这个行为的行政机关"。}你的店铺因工商局开出的罚款决定不服——可以向法院提行政诉讼。在这类程序里：被告\u200c不是某一个具体的自然人——\xuehai{法治的古典源头——从亚里士多德到戴西}{
  古希腊的亚里士多德在《政治学》里最早划了一条线："法治优于一人之治。"不是"法律比人聪明"——而是法律不带脾气、没有私心、不会因为昨晚没睡好而把你的申诉顺手丢进碎纸机。但亚里士多德也警告：法律写得再完美也只是一段静态文字——谁来解释它、谁来执行它——如果解释和执行它的那个人自己不受约束，法治就退化成了"用法来治"。
}

\subsection*{法官的自由裁量权：同一条法律会判出不一样的数字吗？}

法条上写："恶意扰乱公共秩序——拘留五到十天。"一人站在电影院散场时大声接三通电话拒不出场（恶意程度低）。另一个人在周末人海高峰的商场里报假警喊有炸弹——商场数层紧急封闭疏散了数万人、二十家商铺丢失了整半天销售额（恶意极高）。两人都撞上了"扰乱公共秩序"这一条——但判同一等天数，你觉得是公平还是不公平？自由裁量给了法官结合具体情节匹配惩罚力度的空间——但两个不同的法官在非常接近的两个案件里用同一法条写出的刑期的确可能不一样。\textbf{裁判尺子的张力从来没有被彻底调和。}

\subsection*{要点总结}
\begin{itemize}
  \item 立法的过程通常需要经过调研、论证、公开征求意见、审议表决等多道程序——目的是防止"拍脑袋立法"
  \item 三种诉讼类型：民事诉讼（私人纠纷）、刑事诉讼（国家公诉犯罪行为）、行政诉讼（公民告行政机关）
  \item 亚里士多德的"法治优于人治"：法律不带脾气、没有私心——但谁来执行法律本身也需要被约束
  \item 法官的自由裁量权是必要的（让惩罚匹配个案情节），但也可能导致类似案件被不同对待
\end{itemize}

\subsection*{讨论题}
\begin{enumerate}
  \item 两个人偷同样金额的钱——一个为母亲买药，一个纯为享乐。你觉得法官应该"同罪同罚"还是"同罪不同罚"？如果考虑动机，会不会有人学会"编一个悲情故事"来减刑？
  \item 行政诉讼中公民告政府——你觉得这种诉讼在现实中公平吗？政府掌握着资源、信息和法律团队，公民能获得真正的"平等武装"吗？
\end{enumerate}

\newpage
\section{\textcolor{orange!80!black}{第五课\quad 民主——多少种玩法？}}

\subsection*{直接民主——从雅典到瑞士}

公元前五世纪的雅典，全体成年男性公民被召集在普尼克斯山上——他们不需要选"代表"，每一项议题（宣战、财政、放逐将领）由\textbf{本人举手或投碎石}直接表决。今天仍在规模化使用\textbf{直接民主}的是瑞士。每年瑞士公民要面对几场全民投票\u2014\u2014要不要补贴买新型战斗机、要不要提高健康保险的最低自付额度、要不要在某座城市郊外建新的工业园区。\textbf{这种民主让公民手上有一把直接插入具体决策的螺丝刀——但不是每一个操作者都有足够时间把整本工程手册读完。}把复杂的财政方案交给普通人看着一张选票描述决定——其结果有时与技术官僚的专家方案产生严重冲突。

\subsection*{代议制——选一个人替你判断}

现代大多数国家用\textbf{代议制}——公民用选票把一批代表送进议会，由他们来审查、修订、投票通过或否决法律。选的代表如果做的不满意\u2004下一个选举期可以被换掉。关键问题来了：\textbf{你怎么保证被选上的这个人会忠实传达你的意见而不是拿你的声音去交换他自己的便利？}代议制把这个张力放在了核心：\textbf{信任+定期换人的机制}是它被接受的根本条件。

\subsection*{选举制度的潜台词——谁的声音能被挂在票数上}

一个国家的议会把这些代表分配到议员席位上——这步程序设计决定了小声音能不能被对话机放大。

\textbf{英国和美国的作法——赢家通吃}（在各自的一人选区中得最高票即当选）。一个政党拿了全国15\%的分散支持但没有在任何单一选区高过其他人——它可能在议会里\textbf{一个席位都没有}。这套制度逼着所有势力压缩成两个大党交替执政。

\textbf{荷兰和以色列的作法——比例代表制}。一个只拿3\%选票的小党\textbf{几乎必定获得3\%的议会席位}。后果是\u200c议会中常年遍布多个小政党（荷兰最近的一次二院大选有十几个政党获得席位），组阁谈判耗时长——但碎片的声音不会被剔除。

\textbf{德国和新西兰的折中方案——混合选举制}：每个选民投两票——第一票按自己的选区选一个人（赢家通吃），第二票投给一个政党名单（按全国比例分配）。这套设计\textbf{同时保证了每一片地理区有自己的直选代表，也不会把全国少数派的票一刀抹掉}。

\xuehai{迪韦尔热定律——选举制度塑造政党数量}{
  法国政治学家莫里斯·迪韦尔热在1950年代总结了一条被反复验证的经验规律：赢家通吃的单选区制倾向于产生两大党——因为小党的分散票源在任何一个单区都没法"冲顶"，选民不想浪费选票、政党不想浪费资源。比例代表制倾向于产生多党——因为只要跨过最低门槛，分散的少数票也能兑换为议席。不是先有两大党才设计了单选区制——是单选区制在长期运行中把政治力量挤压成了两大阵营。
}

\subsection*{中国的人民代表大会制度——代表是怎么把声音传上去的}

中国实行的是\textbf{人民代表大会制度}——人民通过选举产生各级人大代表，由代表组成各级人民代表大会，统一行使国家权力。

拿一个区的运转来举例。一个区人大代表在社区走访调研的时候，反复听到同一个反映——\textbf{周边几个新建小区入住之后，对口的初中学位不够了，不少小学毕业生面临跨学区调剂}。这位代表整理数据、收集居民联署之后，形成了一份书面的\textbf{代表建议}。这份建议按规定提交给区人大常委会的工作机构——转发给区教育局办理。区教育局在收到建议后必须在规定的时限内书面答复办理结果，并由区政府向区人大常委会报告总体办理情况。

如果教育局的答复代表不满意——他在下一次人代会期间可以提请对口工作委员会对同一问题再次交办、跟催、或者约见教育局负责人当面说明。\textbf{人大代表不是"领工资的公务员"——多数是兼职，他们的初始身份可能是教师、医生、个体户或企业管理人员——但给他们留下的一条窄路是：在闭会期间仍然有权向本级人大常委会提出建议。}

除了人大，还有一套并行系统——\textbf{中国人民政治协商会议}（政协）。政协委员来自各民主党派、无党派人士、各人民团体和各界别（科技界、教育界、医卫界、文艺界等），不直接选举产生，由协商推荐。他们的主要功能是\textbf{政治协商、民主监督、参政议政}——在政府制定规划、修改法律或预算的酝酿阶段，政协可以调研并提出书面意见。政协的提案和人大建议不同——\textbf{没有法律强制答复效力，但通常政府相关机构都会予以书面办理并反馈}。

\texbf{人大定决策，政协提建议——两条平行轨道各有各的功能。}

\subsection*{要点总结}
\begin{itemize}
  \item 直接民主（如雅典和瑞士）让公民直接参与决策，但需要公民有足够的信息和判断力
  \item 代议制民主通过选举代表来治理，核心机制是"信任+定期换人"
  \item 选举制度设计影响政治格局：赢家通吃倾向于两党制，比例代表制倾向于多党制（迪韦尔热定律）
  \item 中国的根本政治制度是人民代表大会制度——人民通过选举代表行使权力
  \item 政协系统与人大并行：人大定决策，政协提建议——两条轨道各有不同功能
\end{itemize}

\subsection*{讨论题}
\begin{enumerate}
  \item 全班举手表决决定去一个有台阶的景点——如果程序正确但结果对轮椅同学不公平，那么"程序的正确"能不能为"结果的错误"辩护？你觉得民主除了"多数决"之外，还需要什么配套原则？
  \item 如果课代表不应该替自己不同意的意见传话——那选他当代表的意义是什么？如果应该——那是不是要求代表放弃自己的判断？这个矛盾如何解决？
\end{enumerate}

\newpage
\section{\textcolor{orange!80!black}{第六课\quad 集权与分权——两种治理逻辑}}

\subsection*{中央集权的好处和代价}

全中国所有的校车规格都由教育部设在北京的一个处室画方案——安全带间距、冷暖空调的最低功率、行驶路线上的哪一座窄桥不许通过。好处是——\textbf{最偏远的山村小学收到的校车和一线城市用的是同一套安全标准}。代价是——川西山区的窄路根本走不了处长画的全国统一长路线，但他无权微调——\textbf{中央看不见每一个山谷的局部条件}。

秦朝把全国切成郡县——官员由中央直接任免、不可世袭、数年被调动轮换。集中动员修驰道、调军粮效率惊人。当叛乱在帝国边缘炸开——每一座郡城没有独力调动资源自组织的灵活性，中央无法向每一个邻接单元同时增援——秦在行政专业化拉到极限时\textbf{失去了应急弹性}。

\subsection*{联邦制——贴近但容易拉扯}

联邦制的底层逻辑是：\textbf{外交、国防、货币和跨州贸易\u200c由联邦政府统一处理}；每个州自收一批税，各自管自己的教育、治安、公共交通和医疗。加州想对新能源车出让人发疯的补贴，它不需要等待密西西比州的批准。但如果新冠疫情穿过州界线冲进养老院——每一个州各自下单抢购呼吸机导致整体价涨——\textbf{联邦的反应迟缓实打实地让初始窗口期白白流失}。

\textbf{所有联邦制国家都永远在谈判同一个图纸上的虚线——"从这一块开始是联邦的事，从那一边向外是州自己的事。"}这条线不是刻死的，是被一次次危机重新推过位置的。

\subsection*{中国——单一制框架下的央地拉锯}

法理上\textbf{中国是一个单一制国家：地方各级政府的权力不是自己原本就有的，而是由中央授予的}。但在实际的财政和行政操作中——央地之间的\textbf{谈判从未停止过}。

1994年税改以前，很大比例的地方税收留在地方。地方财力充足时乡镇企业依靠预算外自留的资金快度孵化出来。94年的分税制一改制——把消费税和增值税（最大头的流转税）大头划归中央——地方政府维持的教育、医疗、养老支出等大量公共责任仍落在它的头上，但缺少与之匹配的自有税基。\textbf{土地财政}应运而生：把城市土地的使用权卖给房地产开发商，用交易环节中产生的大笔出让金填平收支间的缺口。

这个逻辑至今仍在滚动——\textbf{它给城市建造了公路、地铁和学校，但它同时也把财政收入绑在了土地价格长期不能下跌的单向风向上。}

\xuehai{分税制、土地财政与"地方竞争"——一个政治经济学的三角}{
  1994年分税制改革之后，地方政府面对着用20-30\%的地方独享税种去支撑超过70\%的地方公共支出责任这样一个结构性缺口。土地出让金——地方政府把城市土地的使用权卖给房地产开发商获得的收入——变成了填上这个缺口的主要来源。这套逻辑在经济高速增长期运转得相当有效：城市化需要盖新房子，盖房子的开发商需要买入地块——土地出让金稳定流入地方财政。

  但它的结构性风险也同样清晰：一旦土地市场的价格预期逆转，地方政府的核心收入支柱就会随之摇晃。\textbf{这不是"某个政策错了"——这是一个制度环节在特定历史阶段内解决了上一个问题，但把自己的脆弱面推到了下一个阶段。}如何逐步用消费税、个人所得税、房地产持有税等更稳定的税基替代一次性土地出让收入——是当代中国央地财政关系中最核心的未完成议题之一。
}

\subsection*{要点总结}
\begin{itemize}
  \item 中央集权的好处是统一标准（最偏远的山村也能用上统一规格的校车），代价是中央看不见每一个地方的局部条件
  \item 联邦制的底层逻辑：不同层级政府分工——全国性事务归联邦，地方性事务归各州
  \item 中国是单一制国家，但央地关系在实践中通过不断谈判调整——1994年分税制是一个转折点
  \item 分税制后地方财政收入缺口催生了"土地财政"——把土地使用权出让金作为重要收入来源
  \item 土地财政为城市基础设施建设提供了资金，但把地方财政绑在了土地价格不能下跌的单行线上
\end{itemize}

\subsection*{讨论题}
\begin{enumerate}
  \item 中央统一标准 vs 地方因地制宜——你觉得教育应该全国统一教材和考试，还是让各省自己决定？统一的好处和坏处各是什么？
  \item 土地财政让城市建起了学校和地铁，但也推高了房价。如果你是城市管理者，你会怎么设计一套不依赖卖地的公共服务筹资模式？
\end{enumerate}

\newpage
\section{\textcolor{orange!80!black}{第七课\quad 国际法与国际组织——地球上没有"地球警察"}}

\subsection*{联合国安理会——五个常任理事国的否决权}

1945年起草联合国宪章的人把国际联盟（1920—1946）的教训啃了又啃之后，设计了一个极其清晰的结构：\textbf{给战后五个最大的军事强国——美、苏\slash 俄、英、法、中——各一张可以单人否决整个安理会决议的否决票。}逻辑粗暴但真实：国联不包含美国的全部时期、它无法执行任何有武力支持的制裁——最终在日本入侵满洲（1931）、意大利吞并埃塞俄比亚（1935）面前除了反复谴责\u2004全部失效。\textbf{否决权不是五个国家从天上捡来的恩惠——是创始人认为要把最强的武装全部都扣在桌面上才能防止爆发另一场世界大战。}

代价当然极大。当任何一个常任理事国本身卷入冲突或其盟国卷入\u2014\u2014安理会直接瘫痪。\textbf{否决权让最需要联合国的那场冲突反而无法通过安理会授权。}

\subsection*{国际刑事法院——可以判，不能抓}

海牙的\textbf{国际刑事法院}（ICC）起诉的罪名限于\textbf{种族灭绝、危害人类罪、战争罪和侵略罪}——针对的是个人而\u2014\u2014不是国家。但ICC不设自己的警察。它的逮捕令的执行必须依赖\textbf{签约国自愿配合在其境内拘押并将嫌疑人转移到海牙}。当被通缉的嫌疑人碰巧是一个主权国家的现任总统\u2014\u2014他拒绝签批把自己的名字登机的是他自己。ICC的拘押请求在他出国正式踏上另一个签约国的领土之前\u2014\u2014经常被温和地免疫。

\subsection*{WTO——你为什么能看到进口香蕉摆在你小区楼下}

世界贸易组织的争端解决机制曾经被看作全球经济的一颗耀眼之星——\textbf{如果A国对B国的香蕉无理征收高关税，WTO可以授权胜诉方用对等金额的对A国产品加征关税}——报复性关税是WTO能用的最锋利的罚款单。但这种机制对大国和小国的作用是\textbf{不对称的}。小国对大国征收报复关税\u2014\u2014大国的消费者几乎不会感到痛。大国的报复反过来可以让小国的某一个支柱产业瞬间订单腰斩。最近几年当最大的经济体之一持续阻止上诉机构的新法官提名时——\textbf{WTO最重要的终审裁决功能被实际冻结了}\u2014\u2014相当于你在比赛两方都同意遵守同一规则本，但缺少裁判来吹终场哨。

\subsection*{巴黎协定——一份全靠承诺的气候合同}

2015年，195个国家在巴黎签署了一份和此前绝大多数国际条约都不太一样的文件——《\textbf{巴黎协定}》。它没有给任何一个国家设置强制性的减排罚款。它的运转机制是\textbf{各缔约国自己提出"国家自主贡献"}——每个国家自己承诺到某一年把碳排放降到什么水平。协定每隔五年做一次全球盘点，看看全世界的承诺加起来到底够不够把全球升温控制在2℃以内。"约束"不在罚款单上——约束靠的是\textbf{定期点名}：每五年你的排放数据摊在所有国家面前，做不到的承诺会被全球媒体和选民不停地质问。

2017年美国宣布退出，2021年重新加入——同一个气候协定被两个相反方向的总统令来回拉扯。\textbf{当全球最大的历史累计排放国可以因为一次大选的换阁就让国际承诺反复调头——协定里"互信"这层纸的韧性真的够吗？}

\subsection*{联合国维和——一顶蓝盔能做什么、不能做什么}

联合国维和部队的士兵——他们的蓝盔和白色UN标记——被部署在世界各地冲突区的隔离线上。他们的权限严格限制在\textbf{自卫和保卫受权执行的任务}中——不能强行解除武装集团的武器，不能越过授权清单直接干预他国政府军。在卢旺达种族灭绝（1994年）最最危急的那几周，地面上的维和部队规模被大幅削减——因为\u2004安理会无法得出增兵共识。\:\texbf{蓝盔能站在村民和武装分子之间当一道"活人铁丝网"——但当铁丝网背后没有足够的武力后盾和一致的政治意愿撑住它的时候——\x{2004}铁丝网就是一层布。}

\subsection*{主权和全球治理之间永恒的张力}

世界卫生组织面对跨境新发传染病的通报、疫苗接种协作——\textbf{它的全部权限都是建议}。它可以宣布"国际关注的突发公共卫生事件"，警告全球航班应加强筛查——但它不能强制任何主权国家封锁机场、不能强制它公开感染数据、不能派检查员硬闯任何一个国家的P4实验室。

\textbf{全世界绑在同一架喷气式客机的舱壁上——但驾驶操作分散在将近两百个彼此承认"你不可干涉我内政"的独立驾驶舱中。}\:\textbf{全球挑战}\u2014\u2014大流行病、气候变迁、金融危机的跨境传导\u2014\u2014全部都是\textbf{没有地球管理局可以授权统一调度}的超级变量。

\xuehai{国际法的执行困境——为什么"国家主权"是国际法永远翻不过去的墙}{
  国际法和国内法之间最根本的差别不在条文上——在\textbf{执法}上。国内法有三根柱子：警察（抓你）、法院（判你）、监狱（关你）。国际法一根都没有。安理会可以授权制裁和维和行动——但五常的否决权让它在任何一个常任理事国卷入的冲突里自动报废。ICC可以签发逮捕令——但它的法警是签约国自己的警察部队，而签约国可以拒绝配合。

  这不意味着国际法"没用"。它提供了一个共同的语言框架——"\textbf{侵略}"这个词的现代共识、种族灭绝的国际司法定义、贸易争端的仲裁程序——这些规范本身在循环地推动主权国家的行为。但国际法永远做不到和国内法一样硬——\textbf{因为国际法的被执行方和执法方是同一种东西：主权国家}。
}——增加非常任国？限制否决门槛？但当一架巡航导弹从常任国境内发射时\u2014\u2014联合国的纸面裁决拿什么来让发射者按下撤弹钮？
  \item 国际刑事法院签发对某现职国家元首的逮捕令——他的政府拒绝签字移交。ICJ的判决怎么落地？
  \item 如果全球排放碳需要一个统一的碳配额监管体系\u2014\u2014谁来当这个能够处罚大排放国的"执行员"？
\end{enumerate}

\subsection*{要点总结}
\begin{itemize}
  \item 联合国安理会五个常任理事国的否决权是一把双刃剑——它为大国提供了留在体系内的动力，但也让安理会在涉及大国利益的冲突中陷入瘫痪
  \item 国际刑事法院（ICC）可以签发逮捕令，但没有自己的警察部队——执行依赖签约国的自愿配合
  \item 国际法缺少国内法的三根支柱（警察、法院、监狱），其执行依赖主权国家的合作
  \item 全球性挑战（气候变化、大流行病、金融危机）要求全球协作，但主权国家体系使统一行动极其困难
  \item 国际法的真正力量不在于"强制力"，而在于提供了一套共同的语言和规范框架
\end{itemize}

\subsection*{讨论题}
\begin{enumerate}
  \item 常任理事国的否决权让联合国在处理大国冲突时基本无效——那么，是应该取消否决权让联合国更"公正"（但大国可能退出），还是保留否决权让联合国至少"存在"？你觉得哪个选择更现实？
  \item 全球气候变暖需要所有国家一起减排——但每个国家都想"别人多减、我少减"。如果没有一个全球政府来强制执行，你觉得用什么方法能让各国真正合作？
\end{enumerate}

\section*{讨论提示}

\subsection*{第一课（社会契约）讨论题提示}
\begin{itemize}
  \item 关于"三种契约精神的偏好"：可以思考——霍布斯的方案最安全但可能失去自由，洛克的方案保护自由但需要成熟的法治传统，卢梭的方案最民主但对公民素质要求极高。现实中，大多数现代国家是三者混合的产物。
  \item 关于"没有政府的自由"：可以思考——索马里的例子说明，没有统一政府时，"自由"可能变成了"强者欺负弱者的自由"。真正的自由需要一个保护弱者的第三方——而这个第三方本身也可能成为压迫者。这个悖论就是政治哲学的核心难题。
\end{itemize}

\subsection*{第二课（权力合法性）讨论题提示}
\begin{itemize}
  \item 关于"魅力型独裁者 vs 无趣的制度"：可以思考——魅力型领袖能在危机时刻快速决策、凝聚人心，但接班问题无法解决。制度的优势不在"效率"而在"可预期"和"可持续"。历史上两者结合的例子不多但存在——华盛顿用个人魅力奠定了美国总统制的制度化基础。
  \item 关于"输掉选举不认输"：可以思考——民主的运转需要参与者接受一个前提："我有可能输，输了就得让位。"这个前提不是天然的——它需要各方对程序和对手有基本的信任。当信任破裂时，再精巧的制度也可能失效。
\end{itemize}

\subsection*{第三课（宪法）讨论题提示}
\begin{itemize}
  \item 关于"成文宪法 vs 传统"：可以思考——成文宪法也有被架空的可能（比如统治者宣布紧急状态暂停宪法）。传统则可以被打破。最可靠的可能是一者互相补充：成文宪法提供了明文底线，政治传统提供了不成文的自我约束。
  \item 关于"谁来监督法官"：可以思考——这是个"谁来监督监督者"的经典难题。答案通常是"分权和制衡"：法官虽然能审查立法，但法官的任命来自其他权力分支，法官的判决也需要行政分支执行。没有人是最终的监督者——制度通过互相制衡来运作。
\end{itemize}

\subsection*{第四课（法律如何工作）讨论题提示}
\begin{itemize}
  \item 关于"同罪同罚 vs 同罪不同罚"：可以思考——完全统一量刑可以防止偏见，但无法适应个案的千差万别。自由裁量权可以更精准地匹配惩罚与过错，但也可能引入法官的个人偏见。每一套法律体系都在"一致性"和"灵活性"之间寻找平衡。
  \item 关于"公民告政府"：可以思考——政府在诉讼中确实拥有更多资源，但行政诉讼制度的初衷就是提供一个"相对公平"的场域。实践中，律师援助、公益诉讼、媒体关注都可以在一定程度上平衡力量差距。但完全平等武装在理论上是很困难的——这正是为什么"法治"要求政府首先自我约束。
\end{itemize}

\subsection*{第五课（民主）讨论题提示}
\begin{itemize}
  \item 关于"程序正确但结果不公"：可以思考——多数决是民主的程序核心，但它需要与"少数人的基本权利不受多数侵犯"这个原则配合。民主不能只是"数人头"——它还需要宪法保护（如平等权、不歧视条款）来防止程序被用于压迫少数。程序正确并不必然保证结果正当，但程序错误几乎一定导致结果不正当。
  \item 关于"代表要不要替自己不赞同的意见传话"：可以思考——这触及了"代表"这个概念的核心矛盾：代表是"传声筒"还是"受托人"？理论上，代表应当运用自己的判断力，而不是机械传达选民意见。但如果代表完全按自己的判断行事，选民的意志如何体现？不同国家的制度对此有不同的回答。
\end{itemize}

\subsection*{第六课（集权与分权）讨论题提示}
\begin{itemize}
  \item 关于"全国统一 vs 地方自主"：可以思考——统一标准保证底线公平（偏远地区也能享受与大城市同样的教材和考试机会），但牺牲了地方适应性和创新空间。一个可行的思路是"中央定标准，地方出方案"——但这又回到了"谁来监督地方方案是否达标"的老问题。
  \item 关于"替代土地财政"：可以思考——成熟的房地产税（按年征收）可以让地方政府有稳定的现金流而不依赖一次性的土地出让。但开征房地产税面临政治阻力——已经买房的人会觉得自己"被收了两次钱"。这可能需要一个长期的过渡方案。
\end{itemize}

\subsection*{第七课（国际法）讨论题提示}
\begin{itemize}
  \item 关于"否决权的保留 vs 取消"：可以思考——否决权反映了国际政治的现实：大国不会接受一个自己说了不算的体系。取消否决权可能导致大国退出联合国，那联合国的权威和代表性反而更弱。问题的关键可能是如何让否决权的使用成本更高（如要求多个常任国一致反对才能否决某些决议）。
  \item 关于"没有全球政府怎么减排"：可以思考——目前最有希望的路径是"自下而上的压力"：国内政治压力（选民要求政府行动）、市场压力（投资者避开高碳企业）、声誉压力（国家不愿被看作是气候行动的破坏者）。巴黎协定的"点名+定期盘点"机制虽然没有强制力，但它在创造一种新的规范：不减排不再是可以被默默接受的行为。
\end{itemize}

\vspace{0.5cm}
\begin{center}
\textcolor{blue!70!black}{\large —— 从利维坦到国际刑事法院——政治和法律的基本骨架你已经走了一遍。每一层秩序都不是从天上落下来的——全是从上一场混乱的废墟中重新拼出来的。——}
\end{center}

\end{document}
